\chapter{Develop and Validate Scientific Code}\label{ch:code}
\section{From a physical transformation to a stage record}

\begin{headline}
  Scientific code exposes the transformation we want to understand. Write its stage
  record before editing: it reveals missing decisions early, and a blank source or oracle is
  a real gap.
\end{headline}

A reader should find the entry point, follow physically named intermediate states, and
locate the selection, weight, formula or convention behind an output. The workflow
applies to substantial scientific code, physics behavior, architecture and explanatory
rewrites; a framework's native conventions stay authoritative for its interfaces
\cite[Scientific coding; Expose the procedure]{skill:coding}.

\paragraph{Worked example.} Reorganize a script that reads an approved nominal input,
applies an approved selection and fills a weighted histogram of an existing observable, so
that it reads clearly while behaving the same. The example fixes no cuts or bin values and
reports no measured result; its scientific choices must come from the owning analysis source.
Table~\ref{tab:stage} is its record.

\begingroup
\small
\begin{longtable}{L{0.27\textwidth}L{0.66\textwidth}}
  \caption{Applying the required stage record to the illustrative histogram rewrite
    \cite[Establish the scientific contract]{skill:coding}.}\label{tab:stage}\\
  \toprule
  Field & Record before editing \\
  \midrule
  \endfirsthead
  \toprule Field & Record before editing \\ \midrule
  \endhead
  Question & Can the existing transformation be made readable while preserving
  its declared behavior? \\
  \rowcolor{ctpSurface0!30}
  Authoritative source & Exact analysis definition and repository section owning
  selection, observable, weighting, and binning. \\
  Owning layer & Repository and module that own analysis combinations; framework
  calibration stays in its framework layer. \\
  \rowcolor{ctpSurface0!30}
  Input and output & Frozen input identity and schema; histogram artifact plus
  event accounting and comparison evidence. \\
  Named stages & Load inputs; select objects/events; compute observable; apply
  weights; fill and persist the result. Use the actual analysis order. \\
  \rowcolor{ctpSurface0!30}
  Invariants and conventions & Units, selection boundaries, object ordering,
  invalid-value treatment, weight signs, bin edges, and flow handling. \\
  Scale and performance & Expected input size; baseline host, CPU allocation,
  I/O location, wall time, and peak memory where performance matters. \\
  \rowcolor{ctpSurface0!30}
  Oracle & Independent small examples and an inspected baseline artifact whose
  definition is understood; declare comparison tolerances. \\
  Acceptance evidence & Stage comparisons, intended-input execution, focused
  invariant checks, required reviews, and measured performance if applicable. \\
  \rowcolor{ctpSurface0!30}
  Non-goals & No new observable, changed cut, changed normalization, or new
  execution framework unless separately agreed. \\
  \bottomrule
\end{longtable}
\endgroup

\contract\ The transformation and its boundary are now explicit. A blank authoritative source or
oracle would mean we do not yet know what agreement would demonstrate.
\gap\ If the existing output is known to contain a bug, preserving it is not a valid
oracle for that behavior: confirm the intended correction and change the expected result.
A refactor and a definition change can share infrastructure, but their acceptance
questions differ \cite[Numerical evidence]{skill:implementation}.

\section{Give each definition one owner}

\begin{headline}
  Trace the output backward and record one implementation owner for each definition, so a
  convenient copy never becomes a second canonical analysis.
\end{headline}

For the example, find where the binning is chosen, the observable calculated, the event
weight assembled and the selection applied. Independent diagnostics may still check the
production result when their role is explicit (Table~\ref{tab:code-owners})
\cite[Ownership and framework composition]{skill:implementation}.

\begin{table}[H]
  \small
  \begin{tabularx}{\textwidth}{L{0.36\textwidth}Y}
    \toprule
    Definition & Owner \\
    \midrule
    Official calibration and framework decorations & The framework layer. \\
    \rowcolor{ctpSurface0!30}
    Combinations of stable objects & The analysis layer. \\
    Mathematics shared by reconstructed and particle-level definitions & A shared helper,
    not copied formulas. The two definitions may differ intentionally, and the helper must not
    erase a meaningful difference between them. \\
    \rowcolor{ctpSurface0!30}
    Host paths & Runtime configuration, kept apart from physics choices where supported. \\
    \bottomrule
  \end{tabularx}
  \caption{Where definitions live \cite[Ownership and framework
      composition]{skill:implementation}.}
  \label{tab:code-owners}
\end{table}

\paragraph{Read in scientific order.} In the example one should see which selected
events enter the observable and which weight reaches the fill. Small pure helpers make an
independently specified operation easy to test; keep a short linear transformation
together when splitting it would hide the order \cite[Expose the procedure]{skill:coding}.

\paragraph{Configuration-driven frameworks.} The native configuration is already the
visible execution plan:
\begin{itemize}
  \item keep its stage composition and physics choices; add thin domain adapters only where
        needed;
  \item use native includes or composition before custom preprocessing;
  \item generate repetitive catalogues only from frozen, validated inputs, with deterministic
        regeneration evidence \cite[Ownership and framework composition]{skill:implementation}.
\end{itemize}

\section{Implement a narrow path and inspect its evidence}

\begin{headline}
  Prove the stage structure on one bounded path first. Passing that slice supports only
  its declared coverage.
\end{headline}

In the example, one approved nominal input and one existing observable establish the
stages before expanding to every sample and systematic. The slice validates nothing about
untouched branches \cites[Establish the scientific contract]{skill:coding}[Validation
  order]{skill:implementation}.

\paragraph{Before changing anything, preserve the baseline:} source revision, input
identity, configuration, command, environment and comparison artifacts. After the
rewrite, run the same contract and compare stages, not just the final picture.
Machine-readable counts and bin contents keep exact evidence; plots show patterns and
unexpected shapes \cites[Performance]{skill:implementation}[Validation
  order]{skill:implementation}.

\paragraph{Iterate:}
\begin{enumerate}
  \item Change one coherent stage or interface inside the owned boundary.
  \item Run the affected focused check and the bounded intended input when the stage needs
        real-artifact validation.
  \item Inspect intermediate accounting and output differences against the declared expectation.
  \item Fix the first unexplained discrepancy, record the cause, and rerun the affected checks
        before expanding downstream.
\end{enumerate}
At a coherent integration checkpoint, run the owning required suites. Repeating unchanged
test runs adds no new scientific argument; broaden or repeat validation for new changes,
failures or open concerns \cite[Numerical evidence]{skill:implementation}.

\section{Choose tests that protect the scientific claim}

\begin{headline}
  For personal scripts and small libraries, acceptance is led by running the intended data
  and inspecting the numbers and plots. Tests guard what a plausible plot could hide.
\end{headline}

Keep a small automated suite for invariants and demonstrated bugs; reusable numerical
helpers benefit from a few independently calculated examples, while short recipes usually
need less machinery \cite[Numerical evidence]{skill:implementation}. For the histogram,
Table~\ref{tab:code-tests} lists errors that survive a visually plausible result. Choose
only the cases the contract needs, calculate their expected behavior independently, and
keep the reasoning with the test \cite[Numerical evidence]{skill:implementation}.

\begin{table}[H]
  \small
  \begin{tabularx}{\textwidth}{L{0.36\textwidth}Y}
    \toprule
    Case & Exposes \\
    \midrule
    An event exactly at a cut boundary & The wrong inequality. \\
    \rowcolor{ctpSurface0!30}
    Objects with equal values & An unintended tie rule. \\
    Signed weights & A mistaken normalization. \\
    \rowcolor{ctpSurface0!30}
    A value at or beyond a bin edge & Wrong flow handling. \\
    \bottomrule
  \end{tabularx}
  \caption{Test cases for the histogram example.}
  \label{tab:code-tests}
\end{table}

\review\ A regression test must fail against the incorrect behavior it guards. Copying
the implementation's formula into the expected value repeats its assumption. Tests that
freeze source wording, import lists, internal call structure or an evolving command
interface need a demonstrated requirement; they do not replace scientific evidence
\cite[Numerical evidence]{skill:implementation}.

\section{Find the first mismatch upstream}

\begin{headline}
  Validate in dependency order and inspect the first applicable stage before relying on
  the next. A final
  histogram can agree through compensating errors.
\end{headline}

\begin{figure}[H]
  \centering
  \begin{tikzpicture}[node distance=4mm,
      stage/.style={check,text width=128mm,align=left}]
    \node[stage] (a) {\textbf{1. Source and event accounting}: same intended source,
      file and event coverage, explicit treatment of missing or repeated inputs.};
    \node[stage,below=of a] (b) {\textbf{2. Schema, units, identity}: compatible branches
      and units; object multiplicities and event/object identity make later comparisons
      meaningful.};
    \node[stage,below=of b] (c) {\textbf{3. Weights and normalization}: inspect factors
      before their product and its total; keep signed-weight conventions.};
    \node[stage,below=of c] (d) {\textbf{4. Raw and weighted cutflows}: the first selection
      stage where counts or weighted sums diverge.};
    \node[stage,below=of d] (e) {\textbf{5. Inclusive distributions}: object and event
      behavior before particular regions.};
    \node[stage,below=of e] (f) {\textbf{6. Regions and observables}: overlaps, yields and
      the measured quantities under the declared definitions.};
    \node[stage,below=of f] (g) {\textbf{7. Systematic coverage}: which variations execute,
      which paths stay absent or untested.};
    \node[stage,below=of g] (h) {\textbf{8. Response and closure}, where relevant:
      matched/missed/fake accounting, response orientation, acceptance, efficiency,
      closure, stress cases.};
    \node[result,text width=128mm,align=left,below=of h] (i) {\textbf{9. Final result and
        covariance}: interpret only after the upstream contract is checked.};
    \foreach \x/\y in {a/b,b/c,c/d,d/e,e/f,f/g,g/h,h/i}
    \draw[flow] (\x.south)--(\y.north);
  \end{tikzpicture}
  \caption{Validation order and the question each stage resolves
    \cite[Validation order]{skill:implementation}.}
  \label{fig:validation}
\end{figure}

\paragraph{Example.} Selected event identities, raw counts and observable values agree,
but the weighted histogram differs. That focuses the next look on weight factors,
normalization and filling conventions; it does not say which is wrong, so compare their
intermediate evidence. If input coverage already differed, that earlier discrepancy comes
first. This illustrates the order; it diagnoses no current run.

Record the cause at its owning layer: a formula in its canonical implementation, a host
path in runtime configuration, an intended convention in the confirmed scientific
definition. Keep the failed comparison as provenance, so a later reader sees why the
expectation changed \cites[Ownership and framework composition; Numerical
  evidence]{skill:implementation} [Source traceability and
  corroboration]{skill:report-production}.

\section{Measure performance and finish with a reading guide}

\begin{headline}
  When performance matters, measure before and after on a declared basis; an unexplained
  failure of the agreed gate blocks acceptance even when the numbers agree.
\end{headline}

\begin{itemize}
  \item \textbf{Declare the basis:} input size, host, CPU allocation, I/O location, command,
        wall time and peak resident memory, and the allowed regression before the rewrite.
        A comparison across different storage or CPU states that limitation.
  \item \textbf{Keep what has measured value:} single-pass, bounded-memory, lazy, vectorized
        or compiled behavior, optimized behind domain-named functions with a clear
        reference implementation.
  \item \textbf{Add caching or concurrency} only for an observed bottleneck or a required
        recovery property \cite[Performance]{skill:implementation}.
\end{itemize}

The handback holds the implementation, validation evidence, limits and a reading guide.
For the example, the guide lets the reader:
\begin{enumerate}
  \item find the entry point and follow the stage outputs;
  \item locate every formula and configuration choice;
  \item tell framework wiring from analysis logic;
  \item run the narrow and the complete checks;
  \item add one observable without duplicating an existing definition.
\end{enumerate}
It includes measured performance where required and states untested coverage \cite[Verify
  and hand back]{skill:coding}.

\review\ Acceptance requires the owning checks, current documentation, and
any applicable performance gate. Physics and numerical conclusions also require
a council review and ultimately the human domain owner
\cites[Verify and hand back]{skill:coding}[Rounds]{skill:council}.
